\documentclass[11pt]{article}
\usepackage{icmlko}
\begin{document}

\title{부호도 크기도 못 믿는다 --- 축 다섯에서 잰 것 여섯 중 넷이
서른여섯에서 뒤집힌다}
\author{월드모델 실험실}
\date{노트 464 · 2026년 8월 2일}
\maketitle

\begin{abstract}
노트 463 이 노트 $441 \sim 443$ 을 정리하면서 \textbf{노트 442 만
미결}로 남겼다 --- 그것도 축소 모형에서 잰 것이라 다시 재야 했다.
36축 챔피언에서 같은 네 팔을 돌렸다.

\begin{center}
\begin{tabular}{lrr}
\hline
평균 표지의 효과 & 축 다섯 (노트 442) & \textbf{36축} \\
\hline
KR 만화 & $-0.1173$ & $\mathbf{-0.0095}$ \\
비게임 앱 & $-0.2404$ & $\mathbf{-0.0030}$ \\
\hline
\end{tabular}
\end{center}

\textbf{노트 442 는 산다} --- 계산되는 표지는 36축에서도 두 짝을 다
떨어뜨린다. \emph{``집 밖에서 원핫의 $0$ 은 무지의 정직한 표시''}라는
문장은 그대로 선다.

\textbf{그런데 크기가 $12$배 · $80$배 줄었다.} 노트 463 의 기계 그대로다
--- 축이 다섯일 때 \texttt{entry\_friction} 이 입력의 큰 몫이라 ``표지가
독 있는 축을 쓰게 만든다''가 크게 물렸는데, 36축에서는 물 데가 작다.
\end{abstract}

\begin{figure}[t]
\centering
\includegraphics[width=\textwidth]{figs/f1.pdf}
\caption{\textbf{가} --- 같은 실험 셋을 축 다섯과 서른여섯에서.
\textbf{나} --- 여섯 측정을 같은 평면에. 대각선 위에 있었다면 축소
모형이 크기까지 맞은 것인데, 하나도 근처에 없고 \textbf{넷은 부호까지
바뀐다}(붉은색).}
\end{figure}

\section{축소 모형은 어느 쪽으로도 안 믿을 만하다}

이 사이클에 축 다섯에서 잰 것 여섯을 서른여섯에서 다시 쟀다.

\begin{center}
\begin{tabular}{lrrl}
\hline
측정 & 축 다섯 & 36축 & \\
\hline
442 · 평균 표지 KR & $-0.1173$ & $-0.0095$ & 부호 유지 \\
442 · 평균 표지 앱 & $-0.2404$ & $-0.0030$ & 부호 유지 \\
443 · 뒤집기$-$가림 KR & $+0.0664$ & $-0.1698$ & \textbf{뒤집힘} \\
443 · 뒤집기$-$가림 앱 & $+0.2182$ & $-0.0174$ & \textbf{뒤집힘} \\
441 · 가림$-$그대로 KR & $+0.2130$ & $-0.0584$ & \textbf{뒤집힘} \\
441 · 가림$-$그대로 앱 & $+0.3136$ & $+0.0435$ & 부호 유지 \\
\hline
\end{tabular}
\end{center}

\textbf{여섯 중 넷이 부호까지 바뀌고, 부호가 남은 둘도 크기가 한 자릿수
어긋난다.} 축소 모형은 \emph{크기를 어림하는 데도} \emph{방향을 보는
데도} 못 쓴다는 뜻이다.

\textbf{왜 하필 이 실험들인가}도 분명하다. 셋 다 \texttt{entry\_friction}
하나를 어떻게 다룰지에 대한 것인데, 그 축이 입력에서 차지하는 몫이
\textbf{$1/5$ 에서 $1/36$ 로} 줄었다. \emph{한 축을 두고 하는 실험은
그 축의 몫에 비례해서 커진다.} 축소 모형이 그 몫을 일곱 배 부풀렸다.

\section*{규약 --- 무엇을 축소해도 되나}

노트 463 이 \emph{``축소 모형에서 잰 효과는 챔피언에서 다시 재기 전까지
미결''}이라는 조항을 만들었다. 이 노트가 \textbf{언제 축소가 특히
위험한지}를 하나 더한다.

\begin{quote}
\textbf{조항} --- \textbf{한 축(또는 몇 축)을 어떻게 다룰지에 대한
실험은 축 집합을 줄여서 재면 안 된다.} 그 실험의 효과 크기는 그 축이
입력에서 차지하는 몫에 비례하므로, 축을 줄이는 순간 \emph{답이 아니라
문제가 바뀐다.}
\end{quote}

용량 손잡이(깊이 · lr · \texttt{it} · K)처럼 \emph{전체에 걸리는} 것은
축소해도 덜 위험하다 --- 이 사이클의 노트 $457 \sim 461$ 은 전부
챔피언 축 집합에서 쟀고 그래서 이 문제를 안 겪었다.

\section*{남은 것}

\begin{itemize}
\item \textbf{\texttt{OOC\_DROP} 은 여전히 미결이다}(노트 463). 두 짝이
갈리는 한 영샷 기본값은 못 정한다. 셋째 짝이 필요하다.
\item \textbf{노트 444 의 천장 · 노트 453 의 의존도는 축소 모형이
아니었다} --- 짝 자기 자료로만 재거나 챔피언에서 쟀으니 이 병에 안
걸린다. 훑어서 확인했다.
\item \textbf{축 몫을 재는 값싼 자} --- 어떤 축이 입력에서 얼마를
차지하는지 미리 알면, 그 축 실험의 크기를 축 집합 사이에서 환산할 수
있다. 노트 462 의 원핫 의존도와 같은 꼴로 만들 수 있나.
\end{itemize}

\end{document}
